\chapter{ПІДСУМОК}

У результаті порівняння МГЕ і МПГЕ отримали наступні висновки:

\begin{enumerate}
\item
  Коли точка колокації належить елементу границі, по якому проводять
  інтегрування, переважають діагональні елементи у матрицях для МГЕ і
  МПГЕ, що сприяє добрій обумовленості матриці при розв'язуванні системи
  лінійних алгебраїчних рівнянь. Дещо нижча (в порівнянні із МГЕ)
  діагональна перевага у МПГЕ не вплинула на якість обчислень.
\item
  При обчисленні поверхневих зусиль в МГЕ з'являються особливості при
  обчисленні інтегралів, які існують в сенсі Коші, що вимагає
  попереднього аналітичного виділення особливостей (головного значення).
  У МПГЕ всі інтеграли розглядаються як інтеграли без особливостей, а це
  дозволяє у разі потреби обмежуватися тільки числовим інтегруванням.
\item
  Програмне забезпечення у випадку МПГЕ -- єдине для інтегрування по
  внутрішніх і приграничних елементах (криволінійних контурах), а у МГЕ
  слід будувати окреме для граничних (по криволінійних відрізках) і
  внутрішніх елементів.
\item
  В МПГЕ є додатковий параметр методу (ширина приграничної смуги), який
  дає додаткову ступінь вільності при роботі з цим методом. У МГЕ такого
  параметра немає. Як ми бачили з наведеного вище прикладу, в задачах
  теорії пружності він відіграє суттєву роль при збільшені точності
  методу.
\end{enumerate}


